\section{Pairing on the hex lattice}
\label{sec:pairing}

A pairing strategy fixes disjoint pairs of cells before play.  Whenever the
attacker occupies one member of a pair, the defender takes the other.  Such a
strategy blocks a target window when that window contains both members of at
least one pair.  The method would reduce a reply choice to one cell, but the
density required by the three hex-lattice axes is too high for length-six
windows.  Length seven is the exact periodic threshold.

For this section only, a \emph{length-$k$ window} is a set of $k$ consecutive
cells on one of the three axes.  A partial matching is a set of unordered cell
pairs in which each cell occurs at most once.  A matching \emph{covers} a
window if both endpoints of one of its pairs lie in that window.

\subsection{The length-six obstruction}

The first theorem counts the matched cells demanded separately by the three
axis directions.  A cell matched along one direction cannot also meet the
demand from either of the other two directions.

\begin{theorem}[No pairing]
\label{thm:T8}
There is no partial matching $M$ of the axial hex lattice such that every
length-six window contains both cells of a pair in $M$.
\end{theorem}

\begin{plainterms}
A position-independent pairing cannot draw Hexo by itself.  Each axis needs
more than two fifths of the cells to be matched along that axis, while a cell
can belong to only one pair.  The three demands exceed the available cells.
The obstruction concerns a global matching; it does not forbid a partial
pairing inside a certified finite region.
\end{plainterms}

\begin{proof}
Discard every pair that is contained in no length-six window.  This removes
all non-axis-aligned pairs and every axis pair whose endpoints have axis
distance greater than five.  The covering property is unchanged.

Fix one physical axis line and identify its cells with the integers.  A pair
with endpoints $p$ and $p+\delta$, where $1\leq\delta\leq5$, lies in exactly
the windows whose start $s$ satisfies
\[
 p+\delta-5\leq s\leq p.
\]
It therefore covers $6-\delta\leq5$ consecutive window starts.  If pairs are
ordered by their left endpoints, the intervals of starts covered by
consecutive pairs cannot have a gap.  Hence
\[
 p_{i+1}-p_i\leq 6-\delta_{i+1}\leq5.
\]
A line segment of $L$ cells contains $L-5$ internal length-six windows and
therefore needs at least $(L-5)/5$ pairs whose endpoints lie in that segment.
Those pairs account for at least $2(L-5)/5$ matched cells along this axis.

Let
\[
 B_R=\{x:\dist(x,0)\leq R\},
 \qquad |B_R|=3R^2+3R+1.
\]
For any one axis, $B_R$ meets $2R+1$ line segments, and their lengths sum to
$|B_R|$.  Summing the preceding lower bound over those segments and then over
the three axes gives
\[
 \text{matched-cell incidences in }B_R
 \ \geq\ \frac65|B_R|-6(2R+1).
\]
Every matched cell contributes to exactly one axis, namely the axis of its
pair.  A partial matching thus supplies at most $|B_R|$ such incidences.  At
$R=20$, the ball has $1261$ cells, whereas the lower bound is
\[
 \frac65(1261)-6(41)=1267.2>1261.
\]
This is a contradiction.
\end{proof}

The same density calculation gives a threshold for any translation-invariant
line system with negligible boundary.  The hypotheses below state the
geometric setting needed by the count; they do not refer to an arbitrary graph
called a grid.

\begin{corollary}[Translation-invariant line grids]
\label{cor:T8.1}
Consider $k$-in-a-row on a point lattice with $g$ pairwise nonparallel,
translation-invariant line directions.  Suppose that every lattice line is
equipped with length-$k$ windows at all integer offsets and that lattice balls
have boundary-to-volume ratio tending to zero.  If a partial matching covers
every length-$k$ window, then
\[
 g\frac{2}{k-1}\leq1,
 \qquad\text{and hence}\qquad k\geq2g+1.
\]
\end{corollary}

\begin{plainterms}
Each direction asymptotically consumes a fraction $2/(k-1)$ of all cells.
Boundary cells cannot repair an excessive total because their proportion
tends to zero.  For the three hex axes the first possible length is seven, so
Hexo's target length six falls one cell below the pairing threshold.
\end{plainterms}

\begin{proof}
Discard pairs that lie in no target window.  On a fixed direction, a pair at
axis distance $\delta\in\{1,\ldots,k-1\}$ is contained in exactly
$k-\delta\leq k-1$ window starts.  In a large finite set, all but a boundary
term of its internal windows must therefore be covered by at least one pair
per $k-1$ window starts.  Each such pair uses two cells.  The matched-cell
density required for that direction is at least $2/(k-1)$ after division by
the volume and passage along a boundary-to-volume sequence.

Pairs assigned to distinct directions have disjoint endpoints because the
pair family is a matching.  Adding the $g$ lower densities gives
$2g/(k-1)\leq1$.  Rearrangement gives $k\geq2g+1$.
\end{proof}

For comparison, four square-grid directions give the necessary bound
$k\geq9$, consistent with the standard pairing construction for nine in a
row.  Three hex-grid directions give $k\geq7$.  The next theorem attains that
bound on the hex lattice.

\subsection{The threshold construction}

The construction uses a quotient with twelve cell orbits.  Six unit pairs
match those orbits, two pairs in each direction.  Their lifts select one unit
edge phase out of every six on every physical axis line.

\begin{theorem}[Threshold pairing at $k=7$]
\label{thm:T8.2}
On the three-axis hex lattice there is a periodic perfect matching such that
every length-seven window contains exactly one matched pair.  It has a period
lattice of index $12$, and no periodic construction has smaller index.
\end{theorem}

\begin{plainterms}
Increasing the target length from six to seven makes a global pairing
possible.  The twelve-cell tile below matches every cell once and places one
pair in every seven-cell window.  The index bound says that no smaller
periodic tile can do the same job.  This theorem locates the threshold; it does
not provide a pairing defense for length-six Hexo.
\end{plainterms}

\begin{proof}
Write the three positive axis steps as
\[
 H=(1,0),\qquad V=(0,1),\qquad D=(1,-1).
\]
Take the period lattice
\[
 \Lambda=\langle(2,2),(0,6)\rangle
 =\{(2m,2m+6n):m,n\in\mathbb Z\}.
\]
Its determinant has absolute value $12$, and
\[
 F=\{0,1\}\times\{0,1,2,3,4,5\}
\]
is a fundamental domain.  Take every simultaneous $\Lambda$-translate of the
following six pairs:
\[
\begin{array}{c@{\qquad}l}
H &(0,0){-}(1,0),\quad (0,1){-}(1,1),\\
V &(0,3){-}(0,4),\quad (1,3){-}(1,4),\\
D &(1,2){-}(2,1),\quad (1,5){-}(2,4).
\end{array}
\tag{9.1}
\]

\begin{figure}[tb]
\centering
\begin{tikzpicture}[scale=.46]
  % Three copies of the fundamental domain under translation by (2,2).
  \HexPatch{-2/-2,-2/-1,-2/0,-2/1,-2/2,-2/3,
    -1/-2,-1/-1,-1/0,-1/1,-1/2,-1/3,
    0/0,0/1,0/2,0/3,0/4,0/5,
    1/0,1/1,1/2,1/3,1/4,1/5,
    2/2,2/3,2/4,2/5,2/6,2/7,
    3/2,3/3,3/4,3/5,3/6,3/7,
    0/-1,2/1,4/3,4/6}
  \HexShadeList[window one]{0/0,0/1,0/2,0/3,0/4,0/5,
    1/0,1/1,1/2,1/3,1/4,1/5}
  % A few translated pairs show continuation beyond the central tile.
  \HexPair[draw=black!62,line width=1.5pt]{-2}{-2}{-1}{-2}
  \HexPair[draw=black!62,line width=1.5pt]{-2}{-1}{-1}{-1}
  \HexPair[draw=black!62,dashed,line width=1.5pt]{-2}{1}{-2}{2}
  \HexPair[draw=black!62,dashed,line width=1.5pt]{-1}{1}{-1}{2}
  \HexPair[draw=black!62,densely dotted,line width=1.5pt]{-1}{0}{0}{-1}
  \HexPair[draw=black!62,densely dotted,line width=1.5pt]{-1}{3}{0}{2}
  \HexPair[pair horizontal]{0}{0}{1}{0}
  \HexPair[pair horizontal]{0}{1}{1}{1}
  \HexPair[pair vertical]{0}{3}{0}{4}
  \HexPair[pair vertical]{1}{3}{1}{4}
  \HexPair[pair diagonal]{1}{2}{2}{1}
  \HexPair[pair diagonal]{1}{5}{2}{4}
  \HexPair[draw=black!62,line width=1.5pt]{2}{2}{3}{2}
  \HexPair[draw=black!62,line width=1.5pt]{2}{3}{3}{3}
  \HexPair[draw=black!62,dashed,line width=1.5pt]{2}{5}{2}{6}
  \HexPair[draw=black!62,dashed,line width=1.5pt]{3}{5}{3}{6}
  \HexPair[draw=black!62,densely dotted,line width=1.5pt]{3}{4}{4}{3}
  \HexPair[draw=black!62,densely dotted,line width=1.5pt]{3}{7}{4}{6}
  \node[font=\scriptsize,align=center] at \HexAt{-.5}{-2.9}
    {$F-(2,2)$};
  \node[font=\scriptsize,align=center] at \HexAt{.5}{5.9}
    {central $F$};
  \node[font=\scriptsize,align=center] at \HexAt{2.5}{7.9}
    {$F+(2,2)$};
  \node[font=\scriptsize,anchor=west] at (13.7,-1.1)
    {\tikz\draw[pair horizontal] (0,0)--(.7,0);\; $H$};
  \node[font=\scriptsize,anchor=west] at (13.7,-1.8)
    {\tikz\draw[pair vertical] (0,0)--(.7,0);\; $V$};
  \node[font=\scriptsize,anchor=west] at (13.7,-2.5)
    {\tikz\draw[pair diagonal] (0,0)--(.7,0);\; $D$};
\end{tikzpicture}
\caption{The index-$12$ pairing tile.  The softly shaded cells are the
fundamental domain $F$.  The six dark pairs are the representatives in
\textup{(9.1)}; lighter pairs are translates by $\pm(2,2)$.  Diagonal
representatives cross the displayed fundamental-domain boundary, as required
by the quotient matching.}
\label{fig:pairing-tile}
\end{figure}

To prove that the lifted pairs form a matching, use the quotient map
\[
 \phi(q,r)=(q\bmod2,\ r-q\bmod6)
 \colon\mathbb Z^2\longrightarrow\mathbb Z_2\times\mathbb Z_6.
\]
Its kernel is $\Lambda$.  The endpoints in (9.1) have the following images:
\[
\begin{array}{c@{\qquad}l}
H &(0,0){-}(1,5),\quad (0,1){-}(1,0),\\
V &(0,3){-}(0,4),\quad (1,2){-}(1,3),\\
D &(1,1){-}(0,5),\quad (1,4){-}(0,2).
\end{array}
\]
These are the twelve distinct elements of
$\mathbb Z_2\times\mathbb Z_6$.  Every lattice cell is therefore incident
with exactly one lifted pair.

The three axis steps map under $\phi$ to
\[
 (1,5),\qquad(0,1),\qquad(1,4).
\]
Each image has order six.  Each axis direction consequently forms two
six-cycles on the twelve quotient cells.  The construction selects one edge
on each of those six cycles.  Equivalently, the horizontal pair starts on the
two line orbits are $(6n,0)$ and $(6n,1)$, the vertical starts are
$(0,3+6n)$ and $(1,3+6n)$, and both diagonal line orbits have starts with
$q\equiv1\pmod6$.  Thus every physical line has one selected unit-edge phase
modulo six.  A length-seven window has six consecutive internal unit edges,
so it contains exactly one selected pair.

It remains to prove periodic minimality.  We first record the equality forced
on any periodic covering.  Pass, if needed, to a period sublattice whose
quotient is injective on every set of diameter at most six.  Let that quotient
have $N$ cell orbits and $P$ pair orbits.  Discard pairs lying in no
length-seven window.  A remaining pair is axis-aligned at distance
$\delta\in\{1,\ldots,6\}$ and lies in $7-\delta\leq6$ window orbits.  There
are $3N$ window orbits.  Counting pair--window incidences gives
\[
 3N\leq\sum_e(7-\delta_e)\leq6P\leq3N.
\tag{9.2}
\]
The last inequality is $2P\leq N$.  Equality holds throughout.  Hence every
cell is matched, every pair is a unit pair, every window contains exactly one
pair, and each axis has $N/6$ pair orbits.

On one physical line, let $x_s\in\{0,1\}$ indicate that the unit edge starting
at $s$ is selected.  Exact coverage gives
\[
 \sum_{i=0}^{5}x_{s+i}=1
 \qquad\text{for every }s.
\]
Subtracting the equation at $s+1$ from the equation at $s$ yields
$x_s=x_{s+6}$.  Exactly one residue class modulo six is selected, and this
indicator has least positive period six.

Return to any period lattice for the same pairing.  The order of each axis
step in its quotient must be divisible by six, since translation by that
order preserves the line indicator.  The quotient index $N$ is therefore
divisible by six.  If $N<12$, then $N=6$.  An abelian group of order six is
cyclic, and its only elements of order six are $g$ and $-g$.  The difference
of two such elements has order at most three.  The third hex-axis step is the
difference of the first two, yet it too must have order six.  This
contradiction proves that index twelve is minimal.
\end{proof}

\subsection{Finite verification of the periodic search}

The quotient proof reduces the periodic search to six phase variables, one
for each of the two line cycles in each of the three directions.  Each
variable has six possible phases.  An exact-cover row consists of one
line-cycle constraint and the two endpoint constraints of the selected edge.
The resulting instance has $36$ rows and $18$ columns: six line-cycle columns
and twelve cell columns.

The mechanical check exhausts all $6^6=46{,}656$ phase assignments.  An MRV
Algorithm~X traversal visits exactly $419$ recursive states and finds $120$
labelled solutions.  A separate direct traversal of all $46{,}656$ phase
tuples also finds $120$, including the construction above.  The exact
structural checks are recorded in \cref{tab:pairing-checks}.

\begin{table}[tb]
\centering
\caption{Exact checks for the index-$12$ pairing.}
\label{tab:pairing-checks}
\begin{tabular}{lr}
\toprule
quantity & result\\
\midrule
pair orbits & $6$\\
quotient cell orbits checked & $12$\\
quotient cell-incidence histogram & $\{1:12\}$\\
quotient start/axis window orbits checked & $36$\\
quotient contained-pair histogram & $\{1:36\}$\\
cells checked in $[-30,30]^2$ & $3721$\\
finite-patch cell-incidence histogram & $\{1:3721\}$\\
finite-patch windows checked & $11163$\\
finite-patch contained-pair histogram & $\{1:11163\}$\\
\bottomrule
\end{tabular}
\end{table}

The quotient checks are exhaustive rather than samples.  Translation by
$\Lambda$ preserves the matching and all three directions, so every cell is
equivalent to one of the twelve quotient cells and every window is equivalent
to one of the $12\cdot3=36$ quotient windows.  The finite patch is a separate
implementation check.  The exact-cover encoding, direct enumeration, quotient
trace, and patch check are documented with the companion proofs
\cite{hexobot-companions}.

The two threshold theorems leave a narrow implementation conclusion.  A
global static pairing is unavailable for Hexo, but partial pairings remain
available inside position-dependent certified regions.  The index-$12$
construction applies to the neighboring length-seven game and shows that the
density inequality in \cref{cor:T8.1} is attained at $g=3$.
